\documentclass[10pt,a4paper]{amsart}
\usepackage[margin=19mm]{geometry}
\usepackage{amsmath,amssymb,mathtools}
\usepackage[scr=rsfs,cal=euler]{mathalfa}
\usepackage{xfrac}
\usepackage[unicode,hidelinks,bookmarks=false]{hyperref}
\newcommand{\ie}{i.e.\ }
\newcommand{\cf}{cf.\ }
\newcommand{\resp}{respectively\ }
\newcommand{\Subsection}{\subsection}
\providecommand{\nobreakdash}{\nobreak\hbox{-}}
\setlength{\emergencystretch}{2em}
\multlinegap0pt
\usepackage{bm}
\usepackage{bbm}
\usepackage{dsfont}
\usepackage{xspace}
\usepackage{cancel}
\usepackage{mathtools}
\usepackage{amsthm}
\makeatletter
\@ifundefined{theorem}{\newtheorem{theorem}{Theorem}}{}
\@ifundefined{proposition}{\newtheorem{proposition}[theorem]{Proposition}}{}
\makeatother
\title[On the cosine similarity and orthogonality between persisten]{On the cosine similarity and orthogonality between persistence diagrams}
\author{Azmeer Nordin, Mohd Salmi Md Noorani, Nurulkamal Masseran, Mohd Sabri Ismail, Nur Firyal Roslan}
\date{Source: arXiv:2504.04361v1 (CC BY 4.0)}
\begin{document}
\maketitle

\noindent\emph{Source and licence.} On the cosine similarity and orthogonality between persistence diagrams. Version arXiv:2504.04361v1 (CC BY 4.0), distributed under the Creative Commons Attribution 4.0 International licence, as recorded in the arXiv licence metadata for this exact version and shown on the abstract page https://arxiv.org/abs/2504.04361v1. Full rights notice: licenses/arxiv-2504.04361-CC-BY-4.0.txt.\par\smallskip

\noindent\textit{Editorial scope.} Counted unit: the Proposition whose source label is \texttt{proposition: orthogonality trivial perfect} in the article (source0.tex lines 471--473, source label \texttt{proposition: orthogonality trivial perfect}), delivered together with the article's own complete original proof of it (source0.tex lines 474--513, one \texttt{proof} environment of 40 lines). No mathematics is written, repaired or re-derived: statement and proof are the article's own text line for line. Required context retained uncounted: proposition: orthogonal no intersection (lines 422--424). Every definition, display and label of those lines is retained unchanged. Omitted from this package and not redistributed: 1--421, 425--470, 514--793. Nothing inside a retained excerpt is omitted or altered: every retained body is delivered exactly as the article's lines are written, so no omission receipt of any kind accompanies this package. Citations: the retained excerpts carry no citation command at all, so no bibliography entry of the article is transcribed and no reference is redistributed. Reference adaptation: none was needed and none was made; every reference inside the delivered unit resolves inside it, so no unresolved reference or citation is produced. Printed slips kept literal: no printed word, symbol, hypothesis or number of the counted statement or of its proof was altered. Layout and class: the article's own class is \texttt{sn-jnl} with packages including enumerate, xfrac, faktor, lineno, nicefrac, float, amssymb, booktabs, multirow, xcolor, graphicx, amsmath, amsfonts, amsthm; the wrapper is the lane's pinned \texttt{amsart} editorial wrapper at 10pt on A4, which restates the theorem-like environments the retained text needs and adds the article's own macro definitions that the retained text needs (none), with extra wrapper packages (bm, bbm, dsfont, xspace, cancel, mathtools). The delivered numbering is the unit's own. Assets: no retained span contains a figure environment, a graphics-inclusion command, a PDF-page inclusion, a table or a TikZ picture; no image file is copied into this package, the package declares no asset dependency and no image file is redistributed with it. Licence: version arXiv:2504.04361v1 of this article is distributed under the Creative Commons Attribution 4.0 International licence, as recorded in the arXiv licence metadata for this exact version and shown on the abstract page https://arxiv.org/abs/2504.04361v1. A full rights notice is delivered at licenses/arxiv-2504.04361-CC-BY-4.0.txt. Characters: every retained excerpt is pure ASCII, so the delivered engine, XeLaTeX with its default Latin Modern text font, reports no missing character.
\begin{proposition}\label{proposition: orthogonal no intersection}
	The persistence diagrams $D_1$ and $D_2$ are orthogonal if and only if for any pair of points $(b_1, d_1) \in D_1$ and $(b_2, d_2) \in D_2$, the intervals $\left(b_1, d_1\right)_{\mathbb{R}}$ and $\left(b_2, d_2\right)_{\mathbb{R}}$ have an empty intersection.
\end{proposition}

\begin{proposition}\label{proposition: orthogonality trivial perfect}
	If persistence diagrams $D_1$ and $D_2$ are orthogonal, then their trivial matching is perfect under the bottleneck and Wasserstein distances.
\end{proposition}

\begin{proof}
	This clearly holds if either $D_1$ and $D_2$ is empty. We assume that both diagrams are non-empty. For ease of notation, define
	\begin{equation*}
		\omega_{\infty}(\gamma)=\sup _{x \in D_1^{\Delta}}\|x-\gamma(x)\|_{\infty} \quad \textrm{and} \quad \omega_p(\gamma)=\sum_{x \in D_1^{\Delta}}{\|x-\gamma(x)\|_\infty}^p
	\end{equation*}
	for a bijection $\gamma: D_1^{\Delta} \rightarrow D_2^{\Delta}$ in the definitions of bottleneck and Wasserstein distances. Denote $\kappa$ as the trivial matching. For points $x=(b_1, d_1) \in D_1$ and $y=(b_2, d_2) \in D_2$, their paired points on the diagonal line $\Delta$ are
	\begin{equation*}
		\kappa(x)=\left(\frac{b_1+d_1}{2}, \frac{b_1+d_1}{2}\right) \quad \textrm{and} \quad \kappa^{-1}(y)=\left(\frac{b_2+d_2}{2}, \frac{b_2+d_2}{2}\right).
	\end{equation*}
	Note that the points above are the closest to $x$ and $y$ respectively among points on $\Delta$ with respect to supremum norm in $\mathbb{R}^2$.
	
	Suppose that $D_1$ and $D_2$ are orthogonal. We claim that
	\begin{equation}\label{equation: inequality sup norm}
		\sup \big\{\|x-\kappa(x)\|_{\infty},\|y-\kappa^{-1}(y)\|_{\infty}\big\} \leq\|x-y\|_{\infty}
	\end{equation}
	and
	\begin{equation}\label{equation: inequality p-norm}
		{\|x-\kappa(x)\|_\infty}^p+{\|y-\kappa^{-1}(y)\|_\infty}^p \leq {\|x-y\|_\infty}^p.
	\end{equation}
	Proposition \ref{proposition: orthogonal no intersection} states that $\left(b_1, d_1\right)_{\mathbb{R}} \cap\left(b_2, d_2\right)_{\mathbb{R}}=\emptyset$. This implies that either $d_1 \leq b_2$ or $d_2 \leq b_1$. Without loss of generality, we assume the former. The above inequalities are proved respectively as follows:
	\begin{align*}
		\sup \left\{\frac{d_1-b_1}{2}, \frac{d_2-b_2}{2}\right\} & <\sup \left\{d_1-b_1, d_2-b_2\right\} \\
		& \leq \sup \left\{b_2-b_1, d_2-d_1\right\}
	\end{align*}
	and
	\begin{align*}
		2\left(\frac{d_1-b_1}{2}\right)^p+2\left(\frac{d_2-b_2}{2}\right)^p & \leq\left(d_1-b_1\right)^p+\left(d_2-b_2\right)^p \\
		& \leq\left(b_2-b_1\right)^p+\left(d_2-d_1\right)^p\\
		& \leq 2 \cdot \sup\left\{\left(b_2-b_1\right)^p, \left(d_2-d_1\right)^p\right\}.
	\end{align*}
	Consider a bijection $\gamma: D_1^{\Delta} \rightarrow D_2^{\Delta}$. For a point $x \in D_1$, note that if $\gamma(x) \in \Delta$, then
	\begin{equation*}
		\|x-\kappa(x)\|_{\infty} \leq\|x-\gamma(x)\|_{\infty}
	\end{equation*}
	since $\kappa(x)$ is the closest to $x$ among points on $\Delta$. If $\gamma(x)=y \in D_2$, then the inequalities (\ref{equation: inequality sup norm}) and (\ref{equation: inequality p-norm}) hold. These findings conclude that
	\begin{equation*}
		\omega_{\infty}(\kappa) \leq \omega_{\infty}(\gamma) \quad \text { and } \quad \omega_p(\kappa) \leq \omega_p(\gamma)
	\end{equation*}
	for any bijection $\gamma$, and thus $\kappa$ is a perfect matching.
\end{proof}

\end{document}
